\subsubsection{The Eigenspectrum}\label{sec:the-eigenspectrum}

Before connecting QUBO problems to quantum systems, we need to understand one final concept: the \textbf{eigenspectrum} of a Hamiltonian. The definition of eigenvectors and eigenvalues is the following. Given a matrix $A$, an eigenvector $x$ and its corresponding eigenvalue $\lambda$ satisfy the following equation:
\begin{equation*}
    Ax = \lambda x
\end{equation*}
For a Hamiltonian, the same concept appears as (\autopageref{eq:hamiltonian-eigenvalues-eigenvectors}):
\begin{equation*}
    \hat{H} \ket{\psi_{i}} = E_{i} \ket{\psi_{i}}
\end{equation*}
Where $\hat{H}$ is the Hamiltonian matrix, $\ket{\psi_{i}}$ is the $i$-th eigenvector (also called eigenstate), and $E_{i}$ is the corresponding eigenvalue (also called energy level). Therefore, the \hl{possible measurable energies of a quantum system are the eigenvalues of its Hamiltonian.}

\highspace
\textcolor{Green3}{\faIcon{question-circle} \textbf{What does the $\hat{H} \ket{\psi_{i}} = E_{i} \ket{\psi_{i}}$ equation mean?}} It means that when the Hamiltonian acts on the state $\ket{\psi_{i}}$, the result is the same state multiplied by a scalar $E_{i}$. In other words, the \hl{state is not transformed into a different state; it is only scaled}. This is exactly the defining property of an eigenvector. The associated scalar $E_{i}$ is the energy of that state.

\highspace
\textcolor{Green3}{\faIcon{book} \textbf{Stationary States.}} The importance of the eigenvalue equation is not only mathematical. In quantum mechanics, the eigenvectors of the Hamiltonian correspond to \textbf{physically special states}, called \textbf{stationary states}, while the \hl{eigenvalues correspond to the possible measurable energies of the system}. Therefore, studying the eigenspectrum of the Hamiltonian allows us to understand the energy structure of the quantum system.

\highspace
Let's start from a state that is an eigenvector:
\begin{equation*}
    H\ket{v_{1}} = E_{1} \ket{v_{1}}
\end{equation*}
This state has \textbf{one single energy} $E_{1}$. There aren't multiple energies involved. Because there is only one energy, the time evolution is:
\begin{equation*}
    \ket{v_{1}(t)} = e^{-iE_{1}t/\hbar} \ket{v_{1}}
\end{equation*}
The equation above is derived from the Schrödinger equation, which describes how quantum states evolve over time. It is not important to understand the details of the derivation, but the key point is that the whole state gets multiplied by the same phase factor $e^{-iE_{1}t/\hbar}$. Nothing inside the state can ``\emph{get out of sync}'', therefore all measurement probabilities remain unchanged. This is why eigenstates are called \textbf{stationary states}: they do \textbf{not change their measurement probabilities over time}. They evolve only through a global phase factor. Since global phases have no observable effect, all measurement probabilities remain unchanged over time.

\highspace
However, consider a superposition of two eigenstates:
\begin{equation*}
    \ket{\psi} = \frac{1}{\sqrt{2}}\ket{v_1} + \frac{1}{\sqrt{2}}\ket{v_2}
\end{equation*}
Here there are \textbf{two energies} involved $E_{1}$ and $E_{2}$. The time evolution is:
\begin{equation*}
    \ket{\psi(t)} = \frac{1}{\sqrt{2}}e^{-iE_{1}t/\hbar}\ket{v_1} + \frac{1}{\sqrt{2}}e^{-iE_{2}t/\hbar}\ket{v_2}
\end{equation*}
Now the two components of the state evolve with different phase factors. The first component rotates with $E_{1}$, while the second component rotates with $E_{2}$. Then, if $E_1 \neq E_2$, they rotate at different rates. This changes the relative phase between the two components, and that can change measurement probabilities. So the \textbf{state is not stationary}.

\highspace
\begin{flushleft}
    \textcolor{Green3}{\faIcon{question-circle} \textbf{But what does Eigenspectrum mean?}}
\end{flushleft}
The collection of all eigenvalues of a Hamiltonian is called the \definition{Eigenspectrum} of the Hamiltonian. Let's suppose $\hat{H}$ has $n$ eigenvalues $E_1, E_2, \ldots, E_n$. The eigenspectrum is the set $\{E_1, E_2, \ldots, E_n\}$. \hl{Each energy level is associated with one or more eigenstates}.

\highspace
Furthermore, among all eigenvalues $E_1, E_2, \ldots, E_n$, the \textbf{smallest one} is called the \definition{Ground-State Energy}:
\begin{equation}
    E_{\text{ground}} = \min_{i} E_i
\end{equation}
And the correpsonding eigenvector is called the \textbf{ground state} (\autopageref{sec:qubo-as-an-energy-minimization-problem}). Instead, all higher energy levels are called \definition{Excited States}.

\highspace
\textcolor{Green3}{\faIcon{question-circle} \textbf{Why do we care about the Eigenspectrum for a QUBO?}} Recall what we want from a QUBO problem:
\begin{equation*}
    \min E(x)
\end{equation*}
We want the lowest energy configuration. Now look at a quantum system:
\begin{equation*}
    \hat{H} \ket{\psi} = E \ket{\psi}
\end{equation*}
The \textbf{lowest energy is exactly the smallest eigenvalue of the Hamiltonian}. Therefore, if we can encode the QUBO problem into a Hamiltonian, then the solution to the QUBO problem is exactly the ground state of that Hamiltonian. This is the key connection between QUBO problems and quantum systems. By \hl{studying the eigenspectrum of the Hamiltonian, we can find the optimal solution to the original QUBO problem.}